\chapter{RL Training Systems}\label{ch:rlsys}

Chapters~\ref{ch:grpo} and~\ref{ch:grpo2} implemented reinforcement learning with verifiable
rewards (RLVR) as a single Python loop: sample a group of answers, grade them, compute
advantages, take a gradient step, repeat. On one GPU with a 0.6B-parameter model that loop is
easy to read. It is also hopeless at scale. The models that made RL-trained reasoning famous
were trained with thousands of accelerators for weeks, and almost none of that engineering
effort went into the loss function. It went into the \emph{system} that wraps the loss: an
inference engine that produces millions of tokens per step, a training engine that consumes
them, a fleet of reward workers that execute untrusted code, and the plumbing that moves
weights, tokens and log-probabilities between them without leaving expensive hardware idle.

This chapter is about that plumbing. It is written as a graduate systems lecture: we derive cost models, build
a simulator, and study the design space spanned by today's open frameworks (OpenRLHF,
veRL/HybridFlow, TRL, NeMo-Aligner and several newer asynchronous systems). The ecosystem moves
quickly; treat the framework descriptions as a snapshot of designs that were public as of 2025,
and the cost models as the part that will age well.

\begin{figure}[t]
\centering
\resizebox{\linewidth}{!}{%
\begin{tikzpicture}[
  stage/.style={draw=sysblue,fill=sysbluebg,rounded corners=3pt,minimum width=3.0cm,
    minimum height=1.25cm,align=center,font=\small},
  num/.style={circle,fill=sysblue,text=white,font=\bfseries\small,inner sep=1.5pt},
  arr/.style={-{Stealth[length=2.5mm]},thick,sysblue!70}]
\node[stage] (a) at (0,0) {Anatomy of\\an RL step};
\node[stage] (b) at (3.6,0) {Memory \&\\placement};
\node[stage] (c) at (7.2,0) {Weight sync \&\\resharding};
\node[stage] (d) at (10.8,0) {Sync vs.\\async RL};
\node[stage] (e) at (0,-2.0) {Long tail \&\\partial rollouts};
\node[stage] (f) at (3.6,-2.0) {Packing the\\scoring pass};
\node[stage] (g) at (7.2,-2.0) {Train/inference\\mismatch};
\node[stage] (h) at (10.8,-2.0) {Reward \&\\sandboxes};
\foreach \n/\k in {a/1,b/2,c/3,d/4,e/5,f/6,g/7,h/8}
  \node[num] at (\n.north west) {\k};
\draw[arr] (a)--(b); \draw[arr] (b)--(c); \draw[arr] (c)--(d);
\draw[arr] (d.south) to[out=-90,in=90] (e.north);
\draw[arr] (e)--(f); \draw[arr] (f)--(g); \draw[arr] (g)--(h);
\node[draw=labpurple,fill=labpurplebg,rounded corners=3pt,align=center,font=\small,
  minimum width=6.6cm,minimum height=0.9cm] (lab) at (5.4,-3.7)
  {Lab 29: packing, varlen log-probs, TIS, pipeline simulator, sandbox};
\draw[arr,labpurple] (h.south) |- (lab.east);
\end{tikzpicture}}
\caption{Roadmap of this chapter. We start from the dataflow of a single RL step, then add
the systems concerns one at a time, ending with the reward infrastructure. Every stage has
a counterpart in the lab.}
\label{fig:rlsys:roadmap}
\end{figure}

% =====================================================================================
\section{The anatomy of an RL step}\label{sec:rlsys:anatomy}

Strip GRPO (\cref{ch:grpo}) down to its dataflow and five kinds of work remain. They are
summarised in \cref{fig:rlsys:dataflow}.

\begin{enumerate}
\item \textbf{Rollout (generation).} For a batch of $B$ prompts, sample $G$ responses each
  from the current policy $\pi_{\theta}$. This is pure inference: a prefill per prompt
  followed by thousands of memory-bound decode steps (\cref{ch:inference}).
\item \textbf{Reward.} Score each of the $BG$ responses. For RLVR this means extracting an
  answer and checking it with a math verifier (\cref{ch:eval}) or executing generated code
  against unit tests; for RLHF it means a forward pass of a learned reward model.
\item \textbf{Scoring (log-prob) passes.} Compute per-token log-probabilities of every
  response under the policy that generated it ($\piold$) and, if a KL penalty is used,
  under a frozen reference policy $\piref$. These are prefill-only forward passes.
\item \textbf{Advantage computation.} Cheap arithmetic: group means and standard
  deviations (GRPO), or generalised advantage estimation with a critic (PPO).
\item \textbf{Training.} One or more optimizer steps on the clipped surrogate objective,
  each a forward and backward pass over all response tokens.
\end{enumerate}
After the update, the inference engine must receive the new weights before the next rollout
can be on-policy. That arrow---\emph{weight synchronisation}---is what makes RL different
from both pretraining (no inference engine) and serving (no weight updates).

\begin{figure}[t]
\centering
\resizebox{\linewidth}{!}{%
\begin{tikzpicture}[
  box/.style={draw,rounded corners=3pt,minimum height=1.1cm,minimum width=2.7cm,align=center,font=\small},
  gen/.style={box,fill=fmemory!25,draw=napkin},
  rew/.style={box,fill=fdata!25,draw=keygreen},
  trn/.style={box,fill=fcompute!25,draw=sysblue},
  dat/.style={draw=srcgray,dashed,rounded corners=2pt,font=\scriptsize,align=center,inner sep=3pt},
  arr/.style={-{Stealth[length=2.3mm]},thick}]
\node[dat] (p) at (0,0) {prompt\\batch $B$};
\node[gen] (r) at (3.0,0) {Rollout engine\\(vLLM/SGLang)};
\node[rew] (w) at (7.0,1.9) {Reward workers\\(verifier, sandbox,\\reward model)};
\node[trn] (s) at (7.0,-1.9) {Scoring passes\\$\log\piold$, $\log\piref$};
\node[box,fill=white] (a) at (10.8,0) {Advantages\\(group norm)};
\node[trn] (t) at (14.4,0) {Trainer\\(FSDP / Megatron)};
\draw[arr] (p) -- (r);
\draw[arr] (r.north) |- node[pos=0.75,above,font=\scriptsize]{$BG$ texts} (w.west);
\draw[arr] ([xshift=6mm]r.south) |- node[pos=0.72,above,font=\scriptsize]{token ids} (s.west);
\draw[arr] (w.east) -| node[pos=0.25,above,font=\scriptsize]{rewards} (a.north);
\draw[arr] (s.east) -| node[pos=0.25,below,font=\scriptsize]{log-probs} (a.south);
\draw[arr] (a) -- (t);
\draw[arr,fcomm,very thick] (t.south) -- ++(0,-2.9) -| node[pos=0.25,below,font=\scriptsize]
  {weight sync: $\theta_{k+1}$ $\rightarrow$ engine (resharding, broadcast)} ([xshift=-6mm]r.south);
\end{tikzpicture}}
\caption{Dataflow of one RLVR step. Orange = inference workloads (memory-bandwidth bound),
blue = training-style workloads (compute bound), green = reward computation (often CPU-bound,
sometimes another model). The purple edge is the weight synchronisation that closes the loop.}
\label{fig:rlsys:dataflow}
\end{figure}

\subsection{Where does the time go?}

The single most important fact about RL systems for LLMs is that the stages have
\emph{different bottlenecks}. Training and scoring passes process all tokens of a sequence in
parallel and are compute-bound; they want large matrix multiplies and achieve a respectable
fraction of peak FLOP/s (\cref{ch:pretrain}). Generation emits one token per sequence per
step, rereads all weights and the KV cache at every step, and is memory-bandwidth-bound
(\cref{ch:inference}). A configuration that is ideal for one stage is usually wasteful for the
other: training wants the model sharded with FSDP or tensor+pipeline parallelism across many
GPUs and a large activation budget; generation wants as few shards as possible (to avoid
communication on every decode step) and as much free memory as possible for the KV cache.

\begin{napkinmath}[One RL step for an 8B model on 64 H100s]
Take Qwen3-8B ($N\approx 8.19\times10^9$ parameters, 36 layers, 8 KV heads of dimension 128),
$B=512$ prompts, $G=8$ samples each, prompts of 200 tokens and responses of 4{,}000 tokens on
average. That is $BG = 4096$ sequences and $\approx 17.2$M tokens per step.
\begin{itemize}
\item \emph{Training:} $6N\cdot 17.2\text{M} \approx 8.5\times10^{17}$ FLOPs. At 40\% MFU on 64
H100s ($64 \times 989$ TFLOP/s $\times 0.4$) this is about \textbf{33\,s}.
\item \emph{Scoring:} two forward passes ($\piold$ and $\piref$), $2\times 2N\cdot 17.2$M
$\approx 5.6\times10^{17}$ FLOPs, about \textbf{22\,s}.
\item \emph{Generation:} the KV cache costs $2\cdot 36\cdot 8\cdot 128\cdot 2 \approx 147$\,KB per
token. With 16.4\,GB of bf16 weights and $\sim$8\,GB of runtime overhead, about 55\,GB of an
80\,GB card is left for KV; at an average live context of 2{,}200 tokens that is $\approx 170$
concurrent sequences. Each decode step streams $\approx 71$\,GB at 3.35\,TB/s, $\approx$21\,ms,
so one GPU emits $\approx 8{,}000$ tokens/s and 64 GPUs need $\approx$\textbf{32\,s} for 16.4M
generated tokens---\emph{if every slot stayed busy}.
\item \emph{The tail:} if the longest response in the batch is 16{,}000 tokens, the last
sequence decodes alone at $\approx 5$--6\,ms per token (weights dominate), i.e.\ another
\textbf{80--85\,s} during which almost the whole generation fleet is idle.
\end{itemize}
The idealised generation time already matches training, and the tail more than doubles it.
This is why practitioners report that rollout generation typically dominates RL wall-clock
time, and why most of this chapter is about keeping the generator busy.
\end{napkinmath}

\begin{keyidea}[RL is a pipeline of heterogeneous stages]
An RL step chains an inference workload (bandwidth-bound, latency-sensitive, variable length)
with a training workload (compute-bound, fixed shapes) through a weight-update dependency.
Every design in this chapter (colocation, asynchrony, partial rollouts, packing) is an
answer to the question ``how do we keep both kinds of hardware busy despite that dependency?''
\end{keyidea}

% =====================================================================================
\section{Models, memory, and placement}\label{sec:rlsys:placement}

\subsection{Which models must be resident?}

Classic RLHF with PPO juggles four networks: the \emph{actor} (the policy being trained), the
\emph{critic} (a value model, usually initialised from the reward model), a frozen
\emph{reward model}, and a frozen \emph{reference model} used for the KL penalty. GRPO removes
the critic by estimating the baseline from the group of samples; RLVR replaces the reward
model with a program; and recent recipes (DAPO, Dr.~GRPO; \cref{ch:grpo2}) often drop the
reference model by setting the KL coefficient to zero. A modern RLVR run may therefore hold
only the actor---but the actor exists in two \emph{forms}: as training state and as an
inference engine.

\begin{table}[t]
\centering\small
\begin{tabularx}{\linewidth}{@{}>{\raggedright\arraybackslash}Xrr>{\raggedright\arraybackslash}p{3.6cm}@{}}
\toprule
Component & Bytes/param & 8B model & Notes\\
\midrule
Actor training state (bf16 w+g, fp32 master, Adam $m,v$) & 16 & 131\,GB & sharded by FSDP/ZeRO-3\\
Actor activations (per micro-batch) & -- & 10s of GB & checkpointing trades for compute\\
Reference model (bf16, inference only) & 2 & 16.4\,GB & can be offloaded or dropped\\
Critic (PPO only) & 16 & 131\,GB & absent in GRPO\\
Rollout engine weights (bf16) & 2 & 16.4\,GB & per engine replica\\
Rollout KV cache & -- & rest of HBM & sets decode concurrency\\
\bottomrule
\end{tabularx}
\caption{Memory accounting for the models of an RL run (8B example). Training state is
sharded across the trainer's GPUs; inference weights are replicated once per rollout engine
(or sharded by its tensor-parallel degree).}
\label{tab:rlsys:memory}
\end{table}

\subsection{Colocated versus disaggregated placement}

There are two basic ways to place the rollout engine and the trainer on a cluster of $G$
GPUs (\cref{fig:rlsys:placement}).

\paragraph{Colocated (time-shared).} Every GPU hosts both a training shard and an inference
engine, and the two \emph{take turns}. During generation the optimizer state and gradients are
offloaded to host memory (or freed), and the engine's KV cache gets most of HBM; for
training, the engine releases its KV cache (vLLM's ``sleep mode'' and SGLang's memory
release APIs exist for exactly this), and the trainer reclaims memory. veRL's
\emph{3D-HybridEngine} (Sheng et al., HybridFlow, 2025) and OpenRLHF's hybrid-engine mode
follow this pattern. Its virtue is that \emph{no GPU is ever dedicated to a stage that is not
running}. Its costs are the switch overhead (offload/onload, resharding) and the fact that
the two stages are strictly serialised.

\paragraph{Disaggregated (split pools).} $G_{\text{gen}}$ GPUs run inference engines and
$G_{\text{train}} = G - G_{\text{gen}}$ run the trainer. OpenRLHF's original Ray-based design
placed actor, critic, reward and reference models (and vLLM engines) on separate groups of
GPUs; most asynchronous systems use split pools. Each pool can use the parallelism layout
that suits it, and the pools can run \emph{concurrently}---but only if the algorithm tolerates
data generated by slightly older weights (\cref{sec:rlsys:async}). In a strictly synchronous
schedule each pool idles while the other works.

\begin{figure}[t]
\centering
\resizebox{\linewidth}{!}{%
\begin{tikzpicture}[gpu/.style={draw,minimum width=0.95cm,minimum height=0.95cm,font=\scriptsize,align=center},
  lbl/.style={font=\small\bfseries}]
% colocated
\node[lbl] at (2.1,1.5) {(a) Colocated: every GPU time-shares};
\foreach \i in {0,...,3}{
  \node[gpu,fill=fmemory!25] (c\i) at (\i*1.15,0.45) {gen};
  \node[gpu,fill=fcompute!25] at (\i*1.15,-0.55) {train};}
\node[font=\scriptsize,align=left] at (2.1,-1.6) {same HBM, alternating phases;\\switch = offload optimizer, free KV};
% disaggregated
\node[lbl] at (9.8,1.9) {(b) Disaggregated: split pools};
\foreach \i in {0,...,3}
  \node[gpu,fill=fmemory!25] at (6.6+\i*1.15,0) {gen\\TP=1};
\foreach \i in {0,...,3}
  \node[gpu,fill=fcompute!25] (t\i) at (11.8+\i*1.15,0) {train\\FSDP};
\draw[-{Stealth},thick,fcomm] (13.5,-0.6) to[bend left=22] node[below,font=\scriptsize]{NCCL broadcast of $\theta$} (8.3,-0.6);
\draw[-{Stealth},thick,fdata] (8.3,0.6) to[bend left=14] node[above,font=\scriptsize]{rollouts (token ids, log-probs)} (13.5,0.6);
\end{tikzpicture}}
\caption{Two placements. (a) Colocation multiplexes both workloads on the same devices, in
time. (b) Disaggregation dedicates devices to each workload, which permits different
parallel layouts and concurrent execution, at the price of cross-pool weight transfers and
potential idleness.}
\label{fig:rlsys:placement}
\end{figure}

\subsection{A cost model for the two placements}

Let $W_g$ be the generation work of one step and $W_t$ the training+scoring work, both measured
in GPU-seconds at perfect scaling. Ignore the tail for now (we return to it in
\cref{sec:rlsys:tail}). With $G$ GPUs:
\begin{align}
T_{\text{coloc}} &= \frac{W_g}{G} + \frac{W_t}{G} + T_{\text{switch}},\label{eq:rlsys:coloc}\\
T_{\text{sync-split}} &= \frac{W_g}{G_{\text{gen}}} + \frac{W_t}{G - G_{\text{gen}}} + T_{\text{sync}},\label{eq:rlsys:sync}\\
T_{\text{async-split}} &= \max\!\left(\frac{W_g}{G_{\text{gen}}} + T_{\text{sync}},\ \frac{W_t}{G - G_{\text{gen}}}\right).\label{eq:rlsys:async}
\end{align}
Minimising \cref{eq:rlsys:sync} over $G_{\text{gen}}$ gives
$G_{\text{gen}}^\star/G = \sqrt{W_g}/(\sqrt{W_g}+\sqrt{W_t})$ and
$T^\star = (\sqrt{W_g}+\sqrt{W_t})^2/G + T_{\text{sync}} \ge (W_g+W_t)/G$: a synchronous split
is \emph{never} faster than ideal colocation, because each pool idles while the other works.
\Cref{eq:rlsys:async} is minimised by balancing the two terms,
$G_{\text{gen}}^\star/G = W_g/(W_g+W_t)$, which yields $T^\star \approx (W_g+W_t)/G$---the same as
colocation without $T_{\text{switch}}$, but with one step of policy lag. So the trade-off is:

\begin{systemslens}[What disaggregation really buys]
In the idealised model, splitting pools only helps if (i) the algorithm can run the pools
concurrently (asynchrony), or (ii) the stages scale \emph{sub-linearly} on shared hardware so
that specialised layouts recover efficiency (e.g.\ generation with TP$=1$ engines has no
per-token communication, whereas training needs FSDP across many GPUs), or (iii) the switch
cost $T_{\text{switch}}$ is large (offloading 131\,GB of optimizer state over PCIe links that deliver a few tens of GB/s per GPU,
and share host memory bandwidth across 8 GPUs per node, takes seconds per step), or (iv) the hardware is
heterogeneous---older or inference-optimised accelerators can generate while training GPUs
train. In practice, all four effects matter, which is why both designs survive.
\end{systemslens}

% =====================================================================================
\section{Weight synchronisation and resharding}\label{sec:rlsys:sync}

After every optimizer step the inference engines must receive $\theta_{k+1}$. Two problems
hide in that sentence: \emph{moving} the bytes and \emph{reshaping} them.

\subsection{Moving the bytes}

\begin{napkinmath}[How long does a weight update take?]
An 8B model in bf16 is 16.4\,GB.
\begin{itemize}
\item Colocated, same GPU: no transfer at all if the engine can be pointed at the trainer's
tensors (e.g.\ via CUDA IPC handles); otherwise a device-to-device copy at HBM speed, a few ms.
\item Over NVLink within a node ($\sim$450\,GB/s per direction on H100): $\approx 36$\,ms.
\item Over one 400\,Gb/s InfiniBand NIC (50\,GB/s): $\approx 0.33$\,s. With 8 NICs per node
and each trainer rank sending its own shard in parallel, $\approx 0.04$\,s.
\item Over 100\,Gb/s Ethernet (12.5\,GB/s): $\approx 1.3$\,s; for a 32B model, $\approx 5$\,s.
\end{itemize}
A naive point-to-point fan-out to $R$ engines multiplies the sender's time by $R$; a
pipelined (chunked) broadcast tree costs roughly one payload transfer regardless of $R$
(\cref{ch:distributed}). Compared with the tens of seconds of a step these numbers are small
inside a datacenter---but they are \emph{on the critical path} of a synchronous loop, and
they grow quickly for geo-distributed or heterogeneous-cluster RL.
\end{napkinmath}

Two common optimisations reduce the payload: send only the \emph{LoRA adapters} when the
policy is trained with LoRA (\cref{ch:instruct}), so the update is a few hundred MB; and send
\emph{bucketed} chunks so that transfer overlaps with the engine loading the previous chunk.

\subsection{Reshaping the bytes}

The trainer and the engine rarely agree on layout. A trainer using FSDP holds each parameter
flattened and split across $D$ data-parallel ranks; a Megatron trainer splits attention and MLP
weights by tensor-parallel degree $t_{\text{train}}$ and layers by pipeline stage. An inference
engine typically wants whole layers with tensor parallelism $t_{\text{gen}}$, often
$t_{\text{gen}} < t_{\text{train}}$, and fused tensors (e.g.\ a single \texttt{qkv\_proj}
instead of separate $W_Q, W_K, W_V$, or gate and up projections concatenated). Weight sync is
therefore a \emph{resharding} operation: an all-gather (or partial all-gather) along the
training sharding, a re-split along the inference sharding, plus name and layout translation
(\cref{fig:rlsys:reshard}).

\begin{figure}[t]
\centering
\resizebox{0.95\linewidth}{!}{%
\begin{tikzpicture}[sh/.style={draw,minimum height=0.6cm,font=\scriptsize},
  arr/.style={-{Stealth[length=2.2mm]},thick}]
\node[font=\small\bfseries] at (1.6,1.2) {Trainer: FSDP over 4 ranks};
\foreach \i/\c in {0/fcompute!15,1/fcompute!30,2/fcompute!45,3/fcompute!60}
  \node[sh,fill=\c,minimum width=0.8cm] at (\i*0.85+0.3,0.4) {r\i};
\node[font=\scriptsize,align=center] at (1.6,-0.35) {flat shard of every\\parameter on every rank};
\node[draw,rounded corners,fill=fcomm!15,align=center,font=\scriptsize] (g) at (5.6,0.4)
  {all-gather per layer\\+ rename / fuse\\(\texttt{q,k,v}$\to$\texttt{qkv})};
\node[font=\small\bfseries] at (10.0,1.2) {Engine: TP=2 per replica};
\node[sh,fill=fmemory!25,minimum width=1.6cm] (e0) at (9.2,0.4) {cols $0{:}d/2$};
\node[sh,fill=fmemory!45,minimum width=1.6cm] (e1) at (10.9,0.4) {cols $d/2{:}d$};
\node[font=\scriptsize,align=center] at (10.0,-0.35) {column/row splits of\\whole layers (\cref{ch:distributed})};
\draw[arr] (3.6,0.4) -- (g.west);
\draw[arr] (g.east) -- (8.3,0.4);
\end{tikzpicture}}
\caption{Weight synchronisation is a resharding problem: parameters sharded for training
must be gathered, renamed, fused and re-split for the inference engine. veRL's
3D-HybridEngine organises the colocated case so that each GPU only gathers the slice its
inference shard needs, avoiding a full replica in memory.}
\label{fig:rlsys:reshard}
\end{figure}

\begin{pitfall}[Silent layout bugs]
A transposed projection, a swapped \texttt{gate}/\texttt{up} pair, or a RoPE layout mismatch
(interleaved vs.\ rotate-half, \cref{ch:modern}) produces an engine that still emits fluent
text---just from the \emph{wrong} model. Rewards then drift for reasons unrelated to learning.
Always check, after each sync during bring-up, that engine log-probs on a fixed prompt match
the trainer's to within numerical tolerance (\cref{sec:rlsys:mismatch}).
\end{pitfall}
% =====================================================================================
\section{Synchronous versus asynchronous RL}\label{sec:rlsys:async}

\subsection{On-policy in theory, slightly off-policy in practice}

The GRPO objective of \cref{ch:grpo} is an expectation over responses sampled from
$\piold$, the policy at the start of the step. A \emph{synchronous} system enforces this
literally: rollouts for step $k+1$ start only after the weights of step $k$ are in the
engine. An \emph{asynchronous} system lets the generator run ahead. If a sample used in the
update to $\theta_{k+1}$ was produced by $\theta_{k-s}$, we say it has \emph{staleness} $s$.

Off-policy data is not a catastrophe; PPO itself performs several minibatch updates per batch,
so all but the first minibatch is already slightly off-policy, and the clipped ratio
$\pitheta/\piold$ exists precisely to bound how far one batch can push the policy. What
asynchrony changes is that the \emph{behaviour policy} that produced a sample, call it
$\mu = \pi_{\theta_{k-s}}$, is now distinct from the proximal policy around which we clip.
Several systems therefore use a \emph{decoupled} objective (AReaL, Fu et al., 2025, gives a
clear derivation): importance-weight by $\pi_{\text{prox}}/\mu$ to correct for staleness, and
clip $\pitheta/\pi_{\text{prox}}$ to bound the update,
\begin{equation}
\mathcal{L}(\theta) = -\E_{y\sim\mu}\left[\frac{\pi_{\text{prox}}(y)}{\mu(y)}\,
\min\!\Big(r_\theta A,\ \clip(r_\theta,1-\epsilon,1+\epsilon)A\Big)\right],\qquad
r_\theta = \frac{\pitheta(y)}{\pi_{\text{prox}}(y)}.
\label{eq:rlsys:decoupled}
\end{equation}
With $s=0$, $\mu=\pi_{\text{prox}}=\piold$ and \cref{eq:rlsys:decoupled} reduces to the
standard clipped objective. The same structure reappears in \cref{sec:rlsys:mismatch}, where
the behaviour policy differs from $\piold$ not because of staleness but because of numerics.

\subsection{Three schedules}

\Cref{fig:rlsys:timelines} draws the three schedules that bracket the design space.

\begin{figure}[t]
\centering
\resizebox{\linewidth}{!}{%
\begin{tikzpicture}[x=0.55cm,y=0.6cm,
  g/.style={draw=napkin,fill=fmemory!30,minimum height=0.45cm},
  t/.style={draw=sysblue,fill=fcompute!30,minimum height=0.45cm},
  lab/.style={font=\scriptsize\bfseries,anchor=east}]
\newcommand{\blk}[5]{\draw[#1] (#2,#3) rectangle node[font=\tiny]{#5} (#4,#3+0.75);}
% (a) synchronous colocated
\node[lab] at (-0.3,0.4) {(a) sync};
\blk{g}{0}{0}{5}{gen$_0$}\blk{t}{5}{0}{7.5}{train$_0$}
\blk{g}{8}{0}{13}{gen$_1$}\blk{t}{13}{0}{15.5}{train$_1$}
\blk{g}{16}{0}{21}{gen$_2$}\blk{t}{21}{0}{23.5}{train$_2$}
\draw[fcomm,very thick] (7.5,0.375)--(8,0.375); \draw[fcomm,very thick] (15.5,0.375)--(16,0.375);
% (b) one-step async split pools
\node[lab] at (-0.3,-1.35) {(b) async};
\node[lab,font=\tiny] at (-0.3,-1.85) {gen pool};
\node[lab,font=\tiny] at (-0.3,-2.8) {train pool};
\blk{g}{0}{-2.2}{5}{gen$_0$}\blk{g}{5.5}{-2.2}{10.5}{gen$_1$ ($\theta_0$)}
\blk{g}{11}{-2.2}{16}{gen$_2$ ($\theta_1$)}\blk{g}{16.5}{-2.2}{21.5}{gen$_3$ ($\theta_2$)}
\blk{t}{5}{-3.2}{7.5}{train$_0$}\blk{t}{10.5}{-3.2}{13}{train$_1$}\blk{t}{16}{-3.2}{18.5}{train$_2$}
\blk{t}{21.5}{-3.2}{24}{train$_3$}
% (c) partial rollouts
\node[lab] at (-0.3,-4.6) {(c) partial};
\blk{g}{0}{-5}{3}{gen$_0$}\blk{t}{3}{-5}{5.5}{train$_0$}
\blk{g}{6}{-5}{9}{gen$_1$}\blk{t}{9}{-5}{11.5}{train$_1$}
\blk{g}{12}{-5}{15}{gen$_2$}\blk{t}{15}{-5}{17.5}{train$_2$}
\draw[densely dotted,thick,warnred] (3,-4.25) .. controls (4.5,-3.7) .. (6,-4.25);
\draw[densely dotted,thick,warnred] (9,-4.25) .. controls (10.5,-3.7) .. (12,-4.25);
\node[font=\tiny,warnred,align=left,anchor=west] at (17.8,-4.6) {unfinished rollouts\\carried over (mixed $\theta$)};
\draw[-{Stealth},thin] (0,-5.7) -- (24.5,-5.7) node[right,font=\scriptsize]{time};
\end{tikzpicture}}
\caption{Schedules for RL. (a) Synchronous: the generator idles during training (and vice
versa in split pools). (b) One-step asynchronous: batch $k{+}1$ is generated with
$\theta_{k-1}$ or $\theta_k$ while batch $k$ trains; staleness $\le 1$. (c) Partial rollouts:
each generation phase stops once enough responses have finished; long responses continue in
the next phase under newer weights. Purple ticks are weight syncs.}
\label{fig:rlsys:timelines}
\end{figure}

\begin{description}
\item[Synchronous.] Simple and exactly on-policy. Throughput is limited by
$T_{\text{gen}}+T_{\text{train}}$, and the generator's long tail is fully exposed.
\item[One-step (bounded-staleness) asynchronous.] Generation of batch $k+1$ overlaps training
on batch $k$. In steady state the step time is $\max(T_{\text{gen}},T_{\text{train}})$ and the
staleness is bounded by a pipeline depth $s$ (the generator may not run more than $s$ versions
ahead). This is the schedule of several open systems and of our simulator's \texttt{async}
policy. Fully asynchronous variants (AReaL; PipelineRL, which even swaps weights
\emph{in flight} between decode steps of an ongoing sequence) push this further.
\item[Partial rollouts.] Keep more responses in flight than one batch needs; end each
generation phase once $B$ responses have finished and resume the unfinished ones in the next
phase. The tail no longer blocks a step, but some samples are produced by \emph{several}
policies, and the batches over-represent short responses until the long ones catch up. The
Kimi~k1.5 report (2025) describes using partial rollouts for long-context RL.
\end{description}

% =====================================================================================
\section{The long-tail problem}\label{sec:rlsys:tail}

Reasoning-model response lengths are heavy-tailed: most answers to a batch of math problems
end after a few hundred to a few thousand tokens, but a handful run to the length cap. With
continuous batching (\cref{ch:serving}) a decode engine admits new work whenever a slot frees
up, but a \emph{synchronous} RL step has no new work to admit---it needs \emph{these}
$BG$ responses before it can train. The engine's occupancy therefore decays as short
responses finish, and the last few sequences decode almost alone.

We can quantify this with a toy model. Suppose response lengths are log-normal with mean
$\bar L$ and shape $\sigma$, all $n$ responses start together on $n$ slots, and each slot
decodes at a constant $r$ tokens/s. The step's generation time is $\max_i L_i/r$ while the
useful work is $\sum_i L_i/r$, so the generator's utilisation is
\begin{equation}
U = \frac{\sum_i L_i}{n\,\max_i L_i} \;\approx\; \frac{\bar L}{\E[\max_i L_i]}.
\end{equation}
For log-normal lengths $\E[\max_i L_i]$ grows roughly like $e^{\mu + \sigma\sqrt{2\ln n}}$, so
utilisation falls quickly with $\sigma$. \Cref{fig:rlsys:occupancy} shows the decay of
occupancy computed by the simulator from the lab (\ident{generation_occupancy}).

\begin{figure}[t]
\centering
\begin{tikzpicture}
\begin{axis}[width=0.85\linewidth,height=5.2cm,xlabel={time (fraction of step generation time)},
  ylabel={busy decode slots},ymin=0,ymax=1.05,xmin=0,xmax=1,grid=major,grid style={gray!20},
  legend style={font=\scriptsize,at={(0.98,0.97)},anchor=north east},
  every axis plot/.append style={thick}]
\addplot[color=keygreen,mark=none] coordinates {
(0.00,1.000) (0.05,0.766) (0.10,0.445) (0.15,0.289) (0.20,0.164) (0.25,0.117) (0.30,0.102) (0.35,0.086) (0.40,0.047) (0.45,0.039) (0.50,0.031) (0.55,0.023) (0.60,0.016) (0.65,0.016) (0.70,0.008) (0.75,0.008) (0.80,0.008) (0.85,0.008) (0.90,0.008) (0.95,0.008) (1.00,0.000)};;
\addlegendentry{$\sigma=1.0$}
\addplot[color=sysblue,mark=none,dashed] coordinates {
(0.00,1.000) (0.05,1.000) (0.10,0.992) (0.15,0.906) (0.20,0.797) (0.25,0.688) (0.30,0.484) (0.35,0.352) (0.40,0.219) (0.45,0.164) (0.50,0.117) (0.55,0.102) (0.60,0.070) (0.65,0.047) (0.70,0.031) (0.75,0.023) (0.80,0.016) (0.85,0.008) (0.90,0.008) (0.95,0.008) (1.00,0.000)};;
\addlegendentry{$\sigma=0.5$}
\addplot[color=napkin,mark=none,dotted] coordinates {
(0.00,1.000) (0.05,1.000) (0.10,1.000) (0.15,1.000) (0.20,1.000) (0.25,1.000) (0.30,1.000) (0.35,0.961) (0.40,0.867) (0.45,0.789) (0.50,0.688) (0.55,0.477) (0.60,0.320) (0.65,0.195) (0.70,0.125) (0.75,0.094) (0.80,0.047) (0.85,0.031) (0.90,0.008) (0.95,0.008) (1.00,0.000)};;
\addlegendentry{$\sigma=0.25$}
\end{axis}
\end{tikzpicture}
\caption{Occupancy of 128 decode slots during one synchronous generation phase with
log-normal response lengths (simulated with \ident{llmsys.rlsys.sim.generation_occupancy}).
The area above each curve is wasted generator capacity: utilisation is 56\%, 33\% and 13\% for
$\sigma=0.25$, $0.5$ and $1$; with $\sigma=1$ almost 87\% of slot-time is idle because a few
responses run far longer than the rest.}
\label{fig:rlsys:occupancy}
\end{figure}

Mitigations, roughly in order of how much they change the algorithm:
\begin{enumerate}
\item \textbf{Length caps and length-aware rewards.} A hard cap bounds $\max_i L_i$; DAPO
(Yu et al., 2025) adds a soft overlong penalty so the policy learns to avoid truncation.
Capping changes the task (truncated responses score zero) and should be monitored as the
``clip ratio'' of truncated samples.
\item \textbf{Bigger batches.} More concurrent responses amortise the tail across more useful
work, until KV memory runs out.
\item \textbf{Over-provisioning with dropping.} Launch more prompts than needed and train on
the first $B$ groups to finish. This is \emph{biased}: it systematically discards long, often
hard, responses. Use with care.
\item \textbf{Partial rollouts.} Resume, do not drop, the stragglers (\cref{fig:rlsys:timelines}c).
\item \textbf{Asynchrony.} Let the next batch's generation fill the idle slots.
\end{enumerate}

\begin{systemslens}[Dynamic sampling interacts with the tail]
DAPO's \emph{dynamic sampling} (\cref{ch:grpo2}) discards groups whose $G$ rewards are all
equal, since their advantages are zero, and keeps sampling until the batch is full of
informative groups. From a systems standpoint this is an over-provisioning policy with a
data-dependent acceptance rate: early in training most groups fail and late in training most
succeed, so the number of rollouts per optimizer step varies severalfold over a run. An
asynchronous generator that streams groups into a buffer handles this gracefully; a
synchronous one has to guess how many prompts to launch.
\end{systemslens}

\subsection{Simulating the schedules}

The lab's simulator (\ident{llmsys/rlsys/sim.py}) implements the four policies of
\cref{fig:rlsys:timelines} on top of an event-driven model of continuous batching:
\code{n\_slots} decode slots at \code{tokens\_per\_sec} each, log-normal lengths, training
time proportional to tokens, and fixed costs for switching and weight sync.
\Cref{lst:rlsys:rungen} shows its core.

\begin{lstlisting}[caption={Continuous-batching generation with an early-stop for partial
rollouts (from \texttt{llmsys/rlsys/sim.py}).},label={lst:rlsys:rungen}]
def run_generation(pending, n_slots, tps, stop_after, version):
    queue = list(pending)
    heap = []                  # (finish_time, seq, rollout, start_time)
    t, busy, seq = 0.0, 0.0, 0
    finished = []

    def admit(now):
        nonlocal seq
        while queue and len(heap) < n_slots:
            r = queue.pop(0)
            if not r.versions or r.versions[-1] != version:
                r.versions.append(version)          # this policy touched r
            heapq.heappush(heap, (now + r.remaining / tps, seq, r, now))
            seq += 1

    admit(0.0)
    while heap:
        t, _, r, start = heapq.heappop(heap)
        busy += t - start
        r.remaining = 0.0
        finished.append(r)
        if stop_after is not None and len(finished) >= stop_after:
            break                                   # partial rollouts: stop early
        admit(t)
    unfinished = []
    for (ft, _, r, start) in heap:                  # account partial progress
        busy += t - start
        r.remaining = max(0.0, (ft - t) * tps)
        unfinished.append(r)
    return t, finished, unfinished + queue, busy
\end{lstlisting}

\Cref{fig:rlsys:simresults} reports the throughput of each policy as the tail gets heavier.
Two lessons stand out. First, one-step asynchrony helps a lot when the stages are balanced and
the tail is light ($\sigma=0.25$: $+64\%$), but almost nothing when generation is dominated by
a few stragglers---overlapping training with a mostly idle generator does not make the
generator busier. Second, partial rollouts attack the tail directly and keep throughput
roughly constant as $\sigma$ grows, at the price of 45--85\% of samples containing tokens from
more than one policy version and of a length-biased batch composition.

\begin{figure}[t]
\centering
\begin{tikzpicture}
\begin{axis}[ybar,width=0.9\linewidth,height=5.6cm,bar width=7pt,
  symbolic x coords={0.25,0.5,1.0,1.5},xtick=data,
  xlabel={log-normal shape $\sigma$ of response lengths},ylabel={samples / s (simulated)},
  ymin=0,legend style={font=\scriptsize,at={(0.5,1.03)},anchor=south,legend columns=4},
  enlarge x limits=0.18,grid=major,grid style={gray!15}]
\addplot[fill=srcgray!50] coordinates {(0.25,2.013) (0.5,1.448) (1.0,0.744) (1.5,0.449)};
\addplot[fill=fmemory!60] coordinates {(0.25,2.045) (0.5,1.465) (1.0,0.748) (1.5,0.451)};
\addplot[fill=fcompute!60] coordinates {(0.25,3.362) (0.5,2.039) (1.0,0.873) (1.5,0.492)};
\addplot[fill=fdata!70] coordinates {(0.25,2.709) (0.5,2.719) (1.0,2.821) (1.5,3.300)};
\legend{colocated,sync split,async ($s{=}1$),partial}
\end{axis}
\end{tikzpicture}
\caption{Throughput of four RL schedules in the lab simulator (128 slots, batch 128, mean
length 1{,}000 tokens, default costs; \texttt{python labs/lab19\_rlsys/bench.py sim}).
These are outputs of a model, not measurements of a real cluster; their value is in the
trends.}
\label{fig:rlsys:simresults}
\end{figure}

\begin{pitfall}[Partial rollouts bias the batch toward short answers]
In each phase, the responses that finish first are the short ones. Early in a run, the
gradient therefore over-weights short responses, and long responses reach the trainer one or
more versions late. If the reward correlates with length (it usually does for reasoning),
this is a systematic bias. Log the length distribution of \emph{trained} samples, not only of
generated ones.
\end{pitfall}

% =====================================================================================
\section{Packing the scoring and training passes}\label{sec:rlsys:packing}

Once rollouts exist, the trainer runs prefill-style passes over them: the $\piold$ and $\piref$
log-prob passes and the forward/backward of the update. Their cost is linear in tokens (plus a
quadratic attention term), so every padding token is wasted work. With heavy-tailed lengths,
padding to the longest sequence in a micro-batch can waste most of the compute
(\cref{fig:rlsys:packing}).

\begin{figure}[t]
\centering
\resizebox{\linewidth}{!}{%
\begin{tikzpicture}[tk/.style={draw=white,minimum width=0.32cm,minimum height=0.32cm,inner sep=0pt},
  x=0.32cm,y=0.36cm]
\node[font=\small\bfseries,anchor=west] at (0,4.0) {(a) padded micro-batch};
\foreach \r/\L/\c in {0/18/fcompute,1/5/fdata,2/9/fmemory,3/3/fcomm}{
  \foreach \i in {1,...,\L} \node[tk,fill=\c!60] at (\i,-\r+2.6) {};
  \pgfmathtruncatemacro{\s}{\L+1}
  \ifnum\L<18 \foreach \i in {\s,...,18} \node[tk,fill=gray!20] at (\i,-\r+2.6) {}; \fi}
\node[font=\scriptsize,anchor=west] at (0,-1.9) {35 real / 72 computed (49\%)};
\node[font=\small\bfseries,anchor=west] at (22,4.0) {(b) packed row + block-diagonal mask};
\foreach \i in {1,...,18} \node[tk,fill=fcompute!60] at (21+\i,2.6) {};
\foreach \i in {1,...,5} \node[tk,fill=fdata!60] at (21+\i,1.6) {};
\foreach \i in {6,...,14} \node[tk,fill=fmemory!60] at (21+\i,1.6) {};
\foreach \i in {15,...,17} \node[tk,fill=fcomm!60] at (21+\i,1.6) {};
\node[tk,fill=gray!20] at (39,1.6) {};
\node[font=\scriptsize,anchor=west] at (22,0.5) {35 real / 36 computed (97\%)};
% mask
\begin{scope}[shift={(44,-1.8)},x=0.2cm,y=0.2cm]
\fill[gray!12] (0,0) rectangle (17,17);
\fill[fdata!60] (0,17) -- (5,12) -- (0,12) -- cycle;
\fill[fmemory!60] (5,12) -- (14,3) -- (5,3) -- cycle;
\fill[fcomm!60] (14,3) -- (17,0) -- (14,0) -- cycle;
\draw (0,0) rectangle (17,17);
\node[font=\scriptsize,align=center] at (8.5,-2.2) {row 2 mask: attend within\\own segment, causally};
\end{scope}
\end{tikzpicture}}
\caption{Padding vs.\ packing for four responses of lengths 18, 5, 9 and 3. Packing
concatenates sequences into rows of a fixed token budget; a block-diagonal causal mask (or a
varlen kernel driven by \texttt{cu\_seqlens}) prevents attention across sequence boundaries,
and position ids restart at 0 for every segment.}
\label{fig:rlsys:packing}
\end{figure}

Three ingredients make packing correct:
\begin{enumerate}
\item \textbf{Bin packing.} Group sequences into rows under a token budget. First-fit
decreasing (FFD) is simple and provably uses at most $\tfrac{11}{9}\mathrm{OPT}+\tfrac{6}{9}$ bins
(the tight bound is due to D\'osa, 2007).
\item \textbf{Segment-aware attention.} Each token may attend only to earlier tokens of its
own sequence. FlashAttention's \texttt{varlen} interface takes cumulative sequence lengths
(\texttt{cu\_seqlens}) and never materialises the mask; we build the mask explicitly.
\item \textbf{Per-segment positions.} RoPE angles (or GPT-2 position embeddings) must restart
at 0 for each sequence; otherwise the second sequence in a row sees shifted positions.
\end{enumerate}

\begin{lstlisting}[caption={Block-diagonal causal mask and a packed log-prob pass
(\texttt{llmsys/rlsys/packing.py}, abridged).},label={lst:rlsys:packing}]
def block_causal_mask(cu_seqlens):
    N = int(cu_seqlens[-1])
    seg = torch.zeros(N, dtype=torch.long)
    for k in range(len(cu_seqlens) - 1):
        seg[int(cu_seqlens[k]):int(cu_seqlens[k + 1])] = k
    same = seg[:, None] == seg[None, :]
    causal = torch.ones(N, N, dtype=torch.bool).tril()
    return ~(same & causal)                     # True = forbidden

def packed_completion_logprobs(model, seqs, prompt_lens, max_tokens):
    out = [0.0] * len(seqs)
    for bin_ in ffd_pack([len(s) for s in seqs], max_tokens):
        ids, pos, cu = build_packed_row([seqs[i] for i in bin_])
        logp = torch.log_softmax(varlen_forward(model, ids, pos, cu)[0].float(), -1)
        for j, i in enumerate(bin_):
            s, e = int(cu[j]), int(cu[j + 1])
            tok_lp = logp[s:e - 1].gather(-1, ids[0, s + 1:e, None]).squeeze(-1)
            out[i] = float(tok_lp[prompt_lens[i] - 1:].sum())
    return out
\end{lstlisting}

The function \ident{varlen_forward} reuses the \emph{weights} of the core
\ident{Qwen3Model} or \ident{GPTModel} but runs its own attention with the packed mask and
per-token RoPE positions. The lab's test checks that every packed log-prob equals the
unpacked \ident{sequence_logprob} of \cref{ch:refine} to within $10^{-3}$. On a CPU with a tiny
model and 64 log-normal-length sequences, the packed pass took 17\,s against 42\,s for
padded batches of 8 in our runs, with a real-token fraction of 96\% vs.\ 39\%.

\begin{pitfall}[Loss normalisation across micro-batches and ranks]
Token-level losses (DAPO, \cref{ch:grpo2}) are defined as a sum over \emph{all} tokens in the
global batch divided by the global token count. If each micro-batch computes its own
token-mean and gradients are then averaged over micro-batches and data-parallel ranks,
micro-batches with few tokens are silently up-weighted, and the result depends on how the
packer happened to group sequences. The fix is to divide each micro-batch's token \emph{sum} by
the global token count (an all-reduce of a scalar before the backward passes). Several
frameworks have had bugs of exactly this kind.
\end{pitfall}

\begin{systemslens}[Dynamic batch size by token budget]
Fixed-size micro-batches of variable-length sequences produce wildly varying memory peaks.
Production trainers instead form micro-batches by a \emph{token budget} (veRL calls this
dynamic batch size; the idea is general), which makes activation memory roughly constant and
lets the budget be tuned to the largest value that fits. With packing, the token budget is
literally the row length.
\end{systemslens}

% =====================================================================================
\section{Train/inference mismatch}\label{sec:rlsys:mismatch}

A subtle problem appeared once RL systems began using high-throughput inference engines: even
with \emph{identical} weights, the engine's token probabilities $\mu(y_t)$ differ from those the
trainer computes, $\piold(y_t)$. Engines use different kernels (fused attention, paged KV,
different GEMM tilings), different reduction orders, batch-size-dependent algorithms,
sometimes FP8 or quantised KV caches, and sampling-time logit processing. All of these are
harmless for serving and all of them change the low bits of the logits. Over thousands of
tokens and a vocabulary of 150k entries, low-bit differences become occasional large ratio
differences on low-probability tokens---exactly the tokens a sampler sometimes picks.

\begin{figure}[t]
\centering
\resizebox{0.9\linewidth}{!}{%
\begin{tikzpicture}[box/.style={draw,rounded corners,align=center,font=\small,minimum height=1cm},
  arr/.style={-{Stealth},thick}]
\node[box,fill=gray!10] (w) at (0,0) {same weights $\theta$};
\node[box,fill=fmemory!25] (e) at (4.2,1.1) {inference engine\\paged KV, fused kernels,\\bf16/FP8, batch-variant};
\node[box,fill=fcompute!25] (t) at (4.2,-1.1) {trainer\\eager / FSDP kernels,\\bf16, fixed shapes};
\node[box] (mu) at (8.6,1.1) {$\mu(y_t \mid y_{<t})$\\sampled \& logged};
\node[box] (pi) at (8.6,-1.1) {$\piold(y_t \mid y_{<t})$\\recomputed};
\draw[arr] (w) -- (e); \draw[arr] (w) -- (t);
\draw[arr] (e) -- (mu); \draw[arr] (t) -- (pi);
\node[draw=warnred,fill=warnredbg,rounded corners,font=\small,align=center] at (12.2,0)
  {$\rho_t=\dfrac{\piold(y_t)}{\mu(y_t)}\neq 1$\\heavy tails on\\rare tokens};
\draw[arr,warnred] (mu.east) -- (10.9,0.35); \draw[arr,warnred] (pi.east) -- (10.9,-0.35);
\end{tikzpicture}}
\caption{Train/inference mismatch: the samples come from the engine's distribution $\mu$,
while the policy-gradient estimator assumes they come from $\piold$.}
\label{fig:rlsys:mismatch}
\end{figure}

Because $y\sim\mu$, the unbiased policy gradient is
\begin{equation}
\nabla_\theta J = \E_{y\sim\mu}\left[\sum_t \frac{\piold(y_t\mid y_{<t})}{\mu(y_t\mid y_{<t})}
\nabla_\theta \min\!\big(r_t A,\ \clip(r_t,1\pm\epsilon) A\big)\right],
\end{equation}
where we have used a per-token approximation of the sequence-level ratio, as PPO does.
Ignoring $\rho_t = \piold/\mu$ (i.e.\ setting it to 1) introduces a bias that has been
reported to destabilise long runs. Since $\rho_t$ is occasionally huge, the naive correction
has high variance. Two practical estimators, both implemented in
\ident{llmsys/rlsys/mismatch.py}:
\begin{itemize}
\item \textbf{Truncated importance sampling (TIS):} $w_t = \min(\rho_t, C)$ with $C\approx 2$.
This was popularised for LLM RL in a 2025 analysis by Yao et al.\ of the mismatch in
vLLM+FSDP systems; it is the same truncation idea as in off-policy actor--critic methods.
\item \textbf{Masked importance sampling (MIS):} $w_t = \mathbf{1}[1/C\le\rho_t\le C]$, which
drops outlier tokens instead of down-weighting them; variants operate at the sequence level.
\end{itemize}
The weights are \emph{detached}---they rescale gradients and are never optimised.

\begin{lstlisting}[caption={Clipped surrogate with truncated importance-sampling correction.},
label={lst:rlsys:tis}]
def is_weights(logp_train, logp_rollout, cap=2.0, mode="truncate"):
    ratio = torch.exp((logp_train - logp_rollout).detach())
    if mode == "truncate":
        return ratio.clamp(max=cap)
    if mode == "mask":
        return ((ratio >= 1.0 / cap) & (ratio <= cap)).to(ratio.dtype)
    return torch.ones_like(ratio)

def clipped_pg_loss(logp_new, logp_old, adv, mask, clip_eps=0.2,
                    logp_rollout=None, tis_cap=None, tis_mode="truncate"):
    adv = adv[:, None] if adv.dim() == 1 else adv
    ratio = torch.exp(logp_new - logp_old.detach())
    surr = torch.minimum(ratio * adv,
                         ratio.clamp(1 - clip_eps, 1 + clip_eps) * adv)
    if logp_rollout is not None and tis_cap is not None:
        surr = surr * is_weights(logp_old, logp_rollout, tis_cap, tis_mode)
    m = mask.to(surr.dtype)
    return -(surr * m).sum() / m.sum().clamp(min=1)
\end{lstlisting}

Correction is one approach; \emph{elimination} is the other. Making inference kernels
\emph{batch-invariant}---so the result for a sequence does not depend on what else is in the
batch---removes a major source of nondeterminism, and with matching kernels in the trainer it
is possible to run bitwise on-policy RL (Thinking Machines Lab, ``Defeating nondeterminism in
LLM inference'', 2025). Others have argued for higher-precision formats (e.g.\ fp16 with its
larger mantissa instead of bf16) to shrink the mismatch. These remedies cost inference
throughput; TIS costs almost nothing. A reasonable default is: always log the mismatch, always
recompute $\piold$ in the trainer rather than trusting engine log-probs for the ratio, and
enable TIS or MIS.

\begin{napkinmath}[How big is the mismatch from precision alone?]
The lab evaluates the same tiny Qwen3 model in fp32 and in bf16 on random token sequences
(\ident{token_logprobs} with \code{dtype=torch.bfloat16}). On our CPU runs the mean
absolute log-prob difference was $\approx 0.09$ nats, the per-token ratio ranged over
$[0.70,\,1.37]$, and the $k_3$ KL estimate was $\approx 7\times10^{-3}$. A random tiny model
has nearly uniform outputs and is not representative of a trained model; run
\texttt{bench.py mismatch --qwen3\_dir} on real weights to see the long tail on rare
tokens that motivates TIS.
\end{napkinmath}

\begin{keyidea}[Monitor these every step]
Mean $|\log\piold-\log\mu|$; the $k_3$ estimate of $\KL(\mu\,\|\,\piold)$; the fraction of
tokens with $\rho_t\notin[0.5,2]$; the maximum $\rho_t$. A sudden increase is often the first
sign of a weight-sync bug, a kernel change, or a numerically unstable model state.
\end{keyidea}
% =====================================================================================
\section{Configuring the rollout engine for RL}\label{sec:rlsys:engine}

A serving engine tuned for chat traffic is not automatically a good rollout engine. Several
settings that are harmless (or desirable) for serving change the \emph{distribution} that RL
assumes it is sampling from, or waste the structure of RL workloads.

\paragraph{Sampling parameters define $\mu$.} The policy-gradient estimator assumes samples
come from the policy whose log-probabilities appear in the ratio. If the engine samples with
temperature $\tau\neq 1$, top-$p$ or top-$k$ (\cref{ch:scaling}), it samples from a
\emph{transformed} distribution. Either sample at $\tau=1$ without truncation, or compute the
training-side log-probabilities under the same transformation---including the truncation
mask: renormalise over exactly the token set that was available at sampling time (one of the
GRPO modifications listed in \cref{ch:grpo2}). Repetition penalties, logit biases and
grammar-constrained decoding all fall under the same rule.

\paragraph{Exploit the shared prompt.} GRPO samples $G$ responses per prompt. An engine that
supports parallel sampling (\texttt{n}$=G$ in vLLM) or prefix caching (\cref{ch:serving})
prefills each prompt once and forks its KV blocks copy-on-write for the $G$ continuations
(\cref{ch:scaling}). For long prompts (agentic tasks with large contexts, or long
system prompts) this removes up to a factor $G$ of prefill work and of KV memory for the
prompt.

\paragraph{Return token ids, not text.} The trainer must score exactly the token sequence the
engine sampled. Decoding to text and re-tokenising can change the sequence (merges across the
prompt/response boundary, normalisation of whitespace or special tokens), which silently
breaks the ratio. Pass token ids end to end; this is sometimes called a
``token-in, token-out'' interface.

\paragraph{CUDA graphs, cache size and preemption.} Decode steps of small models are dominated
by kernel-launch overhead unless captured in CUDA graphs (\cref{ch:inference}); engines
capture graphs for a set of batch sizes at start-up, which costs time and memory each time an
engine is woken in a colocated setup. The KV cache budget determines concurrency; when it is
exhausted the engine preempts sequences (swap or recompute), which for RL does not change the
result but can multiply the tail.

\begin{systemslens}[Rollout engines are throughput servers with no SLOs]
Serving engines trade throughput for latency SLOs (TTFT, inter-token latency). An RL rollout
has no per-request latency target at all: only the time until the \emph{whole batch} (or the
first $B$ groups) is done matters. Scheduling should therefore maximise concurrency and
prioritise requests that are likely to be long (start them first, a longest-processing-time
heuristic) rather than serve requests in arrival order. Length predictors trained on previous
steps' rollouts have been proposed for exactly this purpose.
\end{systemslens}

% =====================================================================================
\section{Multi-turn and agentic rollouts}\label{sec:rlsys:agentic}

The RLVR loop in this book generates one response per prompt. Agentic RL---training models to
use tools, browse, or edit a code repository---turns each rollout into a \emph{conversation with
an environment}: the model emits a tool call, generation pauses, the environment executes the
call (a shell command, a search query, a unit test), and the observation is appended to the
context before generation resumes (\cref{fig:rlsys:agentic}).

\begin{figure}[t]
\centering
\resizebox{\linewidth}{!}{%
\begin{tikzpicture}[x=0.5cm,y=0.55cm,
  g/.style={draw=napkin,fill=fmemory!30},
  e/.style={draw=keygreen,fill=fdata!25},
  lab/.style={font=\scriptsize\bfseries,anchor=east}]
\newcommand{\seg}[5]{\draw[#1] (#2,#3) rectangle node[font=\tiny]{#5} (#4,#3+0.8);}
\node[lab] at (-0.3,0.4) {rollout 1};
\seg{g}{0}{0}{4}{decode}\seg{e}{4}{0}{9}{tool: run tests}\seg{g}{9}{0}{11}{prefill+dec}
\seg{e}{11}{0}{12.5}{env}\seg{g}{12.5}{0}{16}{decode}
\node[lab] at (-0.3,-0.9) {rollout 2};
\seg{g}{0}{-1.3}{2}{dec}\seg{e}{2}{-1.3}{3}{env}\seg{g}{3}{-1.3}{6}{decode}\seg{e}{6}{-1.3}{14}{tool: slow search}
\seg{g}{14}{-1.3}{17}{decode}
\node[lab] at (-0.3,-2.2) {rollout 3};
\seg{g}{0}{-2.6}{7}{decode}\seg{e}{7}{-2.6}{8}{env}\seg{g}{8}{-2.6}{10}{dec}
\draw[-{Stealth},thin] (0,-3.3) -- (18,-3.3) node[right,font=\scriptsize]{time};
\node[font=\scriptsize,align=left,anchor=west] at (18.6,-0.6)
 {GPU slot is free while a\\rollout waits on its tool:\\admit another sequence,\\keep (or offload) its KV};
\end{tikzpicture}}
\caption{Agentic rollouts interleave decoding (orange) with environment steps (green) whose
latency can exceed the decoding time. An engine that blocks a slot while a tool runs wastes
the slot; asynchronous rollout workers release it and re-admit the sequence when the
observation arrives, ideally with its KV cache still resident (prefix caching).}
\label{fig:rlsys:agentic}
\end{figure}

This changes the systems picture in four ways:
\begin{enumerate}
\item \textbf{Rollouts become state machines.} A rollout worker drives the loop
generate $\rightarrow$ parse tool call $\rightarrow$ execute $\rightarrow$ append observation
$\rightarrow$ generate, typically with \texttt{asyncio} over many concurrent episodes so that
environment latency overlaps other episodes' decoding.
\item \textbf{The KV cache must survive pauses.} Re-prefilling the whole context after every
tool call is quadratic in the number of turns; prefix caching (or explicitly pinned sessions)
keeps the conversation's KV resident between turns.
\item \textbf{Loss masks get more complex.} Only tokens the model generated should receive
gradient; tool outputs and environment observations are context, so they are masked out of
the loss exactly like prompt tokens in SFT (\cref{ch:instruct}). The token-in/token-out rule
matters even more because observations are spliced between generated segments.
\item \textbf{Environments are infrastructure.} Each episode may need its own container
(a repository checkout, a browser, a database), with start-up costs of seconds; environment
pools, snapshotting, and failure handling become as important as the GPU schedule.
\end{enumerate}

% =====================================================================================
\section{Reward infrastructure}\label{sec:rlsys:reward}

In RLVR the reward is a program. That sounds cheap until one counts: 4{,}096 responses per
step, each requiring answer extraction, normalisation and a symbolic-equivalence check
(\cref{ch:eval}), or the execution of model-written code against a test suite. Reward
computation is therefore a \emph{service} with its own throughput, latency, and isolation
requirements (\cref{fig:rlsys:reward}).

\begin{figure}[t]
\centering
\resizebox{\linewidth}{!}{%
\begin{tikzpicture}[box/.style={draw,rounded corners,align=center,font=\small,minimum height=1cm},
  arr/.style={-{Stealth},thick}]
\node[box,fill=fmemory!25] (eng) at (0,0) {rollout\\engine};
\node[box,fill=gray!8,minimum width=2.4cm] (q) at (3.4,0) {request queue\\(per sample,\\streaming)};
\node[box,fill=fdata!20] (math) at (7.4,1.6) {math verifier pool\\(CPU, sympy,\\per-call timeout)};
\node[box,fill=fdata!20] (code) at (7.4,0) {code sandboxes\\(micro-VMs / gVisor,\\no network, rlimits)};
\node[box,fill=fdata!20] (rm) at (7.4,-1.6) {reward-model /\\judge server (GPU)};
\node[box,fill=gray!8] (agg) at (11.2,0) {aggregate\\normalise per\\component};
\node[box,fill=fcompute!25] (tr) at (14.2,0) {advantages\\$\rightarrow$ trainer};
\draw[arr] (eng) -- (q);
\draw[arr] (q.east) -- (math.west); \draw[arr] (q.east) -- (code.west); \draw[arr] (q.east) -- (rm.west);
\draw[arr] (math.east) -- (agg.west); \draw[arr] (code.east) -- (agg.west); \draw[arr] (rm.east) -- (agg.west);
\draw[arr] (agg) -- (tr);
\end{tikzpicture}}
\caption{Reward computation as a service. Samples are scored as soon as they finish, so
reward latency hides behind the generation tail. Different reward types run on different
resources; multi-reward training normalises each component before combining them
(\cref{ch:grpo2}).}
\label{fig:rlsys:reward}
\end{figure}

\paragraph{Streaming.} Because responses finish at very different times, the reward of the
first finished response can be computed while the last ones are still decoding. A reward
service that consumes a stream of completions hides almost all of its latency behind the
generation tail.

\paragraph{Timeouts everywhere.} Symbolic equivalence checking can take unbounded time on
adversarial expressions, and generated code can loop forever. Every reward call needs a hard
wall-clock timeout, a defined reward on timeout (usually 0), and a metric counting timeouts.

\paragraph{Isolation.} Executing model-written code is executing untrusted code. A policy
under optimisation pressure will discover anything that raises its reward: deleting or
monkeypatching the tests, reading the expected outputs from disk, exiting with status 0
before the assertions run, or exhausting resources to crash the grader into a default score.
Production sandboxes use containers with seccomp filters, user-space kernels (gVisor),
micro-VMs (Firecracker), or jail tools (nsjail), with no network, read-only base images,
cgroup CPU/memory limits, and a pool of pre-warmed workers to amortise start-up cost.

\begin{lstlisting}[caption={A teaching sandbox: a fresh interpreter, empty environment,
resource limits and a timeout. Not a security boundary.},label={lst:rlsys:sandbox}]
def _limits(cpu_sec, mem_mb):
    def fn():
        resource.setrlimit(resource.RLIMIT_CPU, (cpu_sec, cpu_sec))
        mem = mem_mb * 1024 * 1024
        resource.setrlimit(resource.RLIMIT_AS, (mem, mem))
        resource.setrlimit(resource.RLIMIT_NOFILE, (32, 32))
        resource.setrlimit(resource.RLIMIT_FSIZE, (1 << 20, 1 << 20))
        os.setsid()
    return fn

def run_python_reward(code, tests, timeout=2.0, mem_mb=512):
    program = code + "\n\n" + tests + "\n"
    t0 = time.perf_counter()
    with tempfile.TemporaryDirectory() as tmp:
        path = os.path.join(tmp, "prog.py")
        with open(path, "w") as f:
            f.write(program)
        try:
            p = subprocess.run([sys.executable, "-I", "-S", path], cwd=tmp,
                               env={}, capture_output=True, text=True,
                               timeout=timeout,
                               preexec_fn=_limits(int(timeout) + 1, mem_mb))
        except subprocess.TimeoutExpired:
            return SandboxResult("timeout", 0.0, time.perf_counter() - t0)
    ok = p.returncode == 0
    status = "pass" if ok else ("fail" if "AssertionError" in p.stderr else "error")
    return SandboxResult(status, float(ok), time.perf_counter() - t0, p.stderr[-500:])
\end{lstlisting}

\begin{pitfall}[Reward hacking through the grader]
\Cref{lst:rlsys:sandbox} concatenates the tests after the code, so a response that calls
\code{os.\_exit(0)} at import time ``passes'' every test. Real graders run tests in a
separate process or namespace from the candidate code, check that the expected number of
tests actually ran, and treat any deviation from the protocol as failure. Audit samples with
suspiciously high reward.
\end{pitfall}

\begin{napkinmath}[Sandbox capacity]
A cold Python process costs tens of milliseconds to start; a container, hundreds of
milliseconds; a micro-VM from a snapshot, on the order of 100\,ms. If 4{,}096 code responses
per step each run a 0.5\,s test suite and the generation phase lasts 60\,s, we need at least
$4096\times 0.5/60\approx 34$ concurrent workers just to keep up with the stream---and several
times that to keep reward latency off the critical path when completions arrive in a burst at
the end of a synchronous phase.
\end{napkinmath}

% =====================================================================================
\section{A tour of open RL frameworks}\label{sec:rlsys:frameworks}

With the vocabulary of the previous sections we can place the open frameworks in a common
design space (\cref{tab:rlsys:frameworks}). The table describes designs as they were
documented around 2025; all of these projects evolve quickly, often converging on each
other's best ideas, so consult their current documentation before relying on a detail.

\begin{table}[t]
\centering\small
\begin{tabularx}{\linewidth}{@{}l>{\raggedright\arraybackslash}X>{\raggedright\arraybackslash}X>{\raggedright\arraybackslash}X@{}}
\toprule
Framework & Orchestration & Training / rollout engines & Placement \& notable ideas\\
\midrule
TRL (Hugging Face) & single program on top of \texttt{accelerate} &
Transformers + DeepSpeed/FSDP; vLLM for generation & vLLM either as a separate server
process or colocated in the training process; easiest entry point, HF-ecosystem native\\
OpenRLHF (Hu et al., 2024) & Ray actors & DeepSpeed ZeRO; vLLM &
originally separate GPU groups per model (actor, critic, reward, reference, vLLM); later a
hybrid-engine colocation mode and asynchronous options\\
veRL / HybridFlow (Sheng et al., 2025) & hybrid: single controller (Ray driver) for the
RL dataflow, SPMD workers inside each model & FSDP or Megatron-LM; vLLM or SGLang &
3D-HybridEngine for colocated, low-redundancy resharding between training and generation
layouts; token-budget dynamic batching\\
NeMo-Aligner (Shen et al., 2024) & Megatron-based distributed jobs & Megatron-LM;
TensorRT-LLM-accelerated generation & large-scale PPO/RLHF in the NVIDIA stack; continued
in NVIDIA's newer NeMo-RL\\
AReaL (Fu et al., 2025) & asynchronous generation and training services &
SGLang rollouts; FSDP/Megatron training & fully asynchronous with bounded staleness and a
decoupled PPO objective; interruptible rollouts\\
\bottomrule
\end{tabularx}
\caption{Design choices of several open RL frameworks for LLMs (snapshot of public designs
around 2025). Other actively developed systems include slime, ROLL, SkyRL, PipelineRL and
prime-rl; the space changes month to month.}
\label{tab:rlsys:frameworks}
\end{table}

\subsection{Single controller versus multi-controller}

The HybridFlow paper frames a useful distinction. In a \emph{multi-controller} (SPMD) program,
every GPU runs the same script and coordination is implicit in collectives---the natural style
for training a single model (\cref{ch:distributed}). In a \emph{single-controller} program, one
driver issues remote calls to groups of workers---the natural style for a dataflow of several
models, because the RL algorithm reads like the loop of \cref{ch:grpo}:

\begin{lstlisting}[caption={The RL loop as a single-controller program (pseudocode in the
style of Ray-based frameworks). Each call fans out to an SPMD worker group.},
label={lst:rlsys:controller}]
for step in range(num_steps):
    batch = dataloader.next()
    rollouts = rollout_group.generate(batch, n=G)           # vLLM/SGLang workers
    rewards = reward_service.score(rollouts)                # CPU pool / sandboxes
    old_lp = actor_group.compute_logprobs(rollouts)         # FSDP workers, packed
    ref_lp = ref_group.compute_logprobs(rollouts) if beta > 0 else None
    adv = group_normalise(rewards, G)                       # on the driver
    metrics = actor_group.update(rollouts, adv, old_lp, ref_lp)
    rollout_group.load_weights(actor_group.export_weights())  # weight sync
\end{lstlisting}

The driver handles only small metadata (or references to data held by workers, as Ray object
references), so it does not become a bandwidth bottleneck, while each \code{*\_group} call is an
SPMD collective program internally. HybridFlow argues that this hybrid gives the
flexibility of single-controller dataflow with the efficiency of SPMD execution.

% =====================================================================================
\section{Checkpointing and fault tolerance}\label{sec:rlsys:ft}

An RL run's state is larger than a pretraining run's: actor weights and optimizer state,
critic state (PPO), the frozen reference (re-loadable), data-loader position and RNG state,
\emph{and} the in-flight rollout state of an asynchronous pipeline---the replay buffer of
finished but untrained samples, partially generated responses, and the policy version that
produced each. Two pragmatic rules:
\begin{itemize}
\item \textbf{Checkpoint at step boundaries and drop in-flight rollouts on restart.} The lost
work is at most one step of generation, and restart logic stays simple. Record the policy
version in every sample so that resumed async pipelines can enforce the staleness bound.
\item \textbf{Make rollout engines stateless and restartable.} Engines crash (OOMs from
unexpectedly long sequences, driver errors). A controller that can re-launch an engine,
reload the latest weights and re-issue its unfinished prompts turns a crash into a slowdown.
\end{itemize}
Sharded, asynchronous checkpointing of the trainer (\cref{ch:pretrain}) applies unchanged.

% =====================================================================================
\section{Putting it together: a design worksheet}\label{sec:rlsys:design}

\begin{napkinmath}[Sizing an 8B RLVR run on 64 GPUs]
Using the numbers from \cref{sec:rlsys:anatomy}: $W_t \approx 64\times(33+22)$\,GPU-s
$\approx 3{,}500$\,GPU-s of training and scoring per step, and
$W_g \approx 64 \times 32 \approx 2{,}050$\,GPU-s of generation if slots stay busy.
\begin{itemize}
\item \emph{Colocated, synchronous:} $(W_g+W_t)/64 \approx 87$\,s plus the tail ($\sim$80\,s) plus
switch overhead: $\sim$170--180\,s per step.
\item \emph{Split, one-step async:} balance $G_{\text{gen}}/G = W_g/(W_g+W_t)\approx 0.37$,
i.e.\ 24 generation GPUs and 40 training GPUs; the per-step time approaches
$\max(W_g/24, W_t/40)\approx 87$\,s, but the tail now sits inside the generation term: unless it
is mitigated (partial rollouts, caps), generation takes $\approx 85+80 = 165$\,s and the
trainer idles half the time.
\item \emph{Split, async + partial rollouts:} with the tail absorbed into the next phase, step
time approaches the balanced $\approx 90$\,s, a $\sim$2$\times$ improvement over the synchronous
baseline, at the cost of staleness $\le 1$--2 and mixed-policy samples.
\end{itemize}
These figures are idealised, but they show where to look first: the tail, then the balance
between pools, then the constant overheads.
\end{napkinmath}

A checklist for a new RL system:
\begin{enumerate}
\item Measure the response-length distribution before choosing a schedule; it determines
whether asynchrony or tail mitigation matters more.
\item Recompute $\piold$ in the trainer and log mismatch metrics; enable TIS/MIS.
\item Pack the scoring and training passes; normalise token-level losses by the global token
count.
\item Bound staleness explicitly, record policy versions per sample, and correct with
importance weights where staleness is non-zero.
\item Stream rewards; put timeouts and isolation around every reward call.
\item Verify weight sync with a log-prob comparison during bring-up.
\end{enumerate}

% =====================================================================================
\section{Summary}\label{sec:rlsys:summary}
\begin{itemize}
\item An RL step chains bandwidth-bound generation, CPU-bound or GPU-bound rewards, and
compute-bound scoring/training through a weight-update dependency. Generation usually
dominates wall-clock time, mostly because of the long tail of response lengths.
\item Colocation time-shares devices (no idle pools, switch overhead); disaggregation
dedicates devices (specialised layouts, concurrency) but only beats colocation when combined
with asynchrony, heterogeneous hardware, or high switch costs.
\item Weight sync is a resharding problem as much as a bandwidth problem; validate it by
comparing log-probs.
\item Asynchronous RL bounds staleness and corrects for it with importance weights (decoupled
objectives); partial rollouts tame the tail at the cost of mixed-policy, length-biased batches.
\item Packing with block-diagonal masks and per-segment positions removes padding waste in
scoring and training passes; loss normalisation must use global token counts.
\item Inference engines sample from $\mu\neq\piold$ even with identical weights; truncated or
masked importance sampling corrects this cheaply, batch-invariant kernels remove it.
\item Rewards are a service: stream them, time them out, and isolate code execution; assume the
policy will find any loophole in the grader.
\end{itemize}

% =====================================================================================
\section{Exercises}\label{sec:rlsys:exercises}

\begin{exercise}[Optimal split]
Derive the optimal $G_{\text{gen}}$ for \cref{eq:rlsys:sync} and show that
$T^\star_{\text{sync-split}} \ge T_{\text{coloc}} - T_{\text{switch}}$ with equality iff one stage
has zero work. For which values of $T_{\text{switch}}$ does the synchronous split win?
\end{exercise}

\begin{exercise}[KV budget]
Repeat the generation napkin math of \cref{sec:rlsys:anatomy} for Qwen3-4B (36 layers, 8 KV
heads, head dimension 128) on an 80\,GB GPU with average live context 6{,}000 tokens. How many
concurrent sequences fit, and what is the upper bound on decode tokens/s per GPU?
\end{exercise}

\begin{exercise}[Staleness bound]
In the \texttt{async} policy of the simulator, generation of batch $k$ may start only after
training step $k-s-1$ has finished. Prove that every sample of batch $k$ has staleness at most
$s$, and show with a timeline that $s=0$ recovers the synchronous split schedule.
\end{exercise}

\begin{exercise}[Loss normalisation bug]
Two micro-batches contain 100 and 1{,}000 response tokens with mean token losses $\ell_1$ and
$\ell_2$. Compute the effective loss when (a) micro-batch means are averaged, and (b) the
token-level definition is used. Write the one-line fix for a gradient-accumulation loop.
\end{exercise}

\begin{exercise}[TIS bias and variance]
For a single token with $\rho\sim$ a log-normal distribution with $\E_\mu[\rho]=1$, compute the
bias $\E_\mu[\min(\rho,C)]-1$ as a function of $C$ and the log-normal variance. Why is a
small bias acceptable here?
\end{exercise}

\begin{exercise}[Packing efficiency]
Show that FFD never uses more than $2\,\mathrm{OPT}$ bins with a short argument (the full
$\tfrac{11}{9}$ bound is harder). Then use \ident{efficiency_report} to compare padding,
length-sorted padding and packing on log-normal lengths with $\sigma\in\{0.25,0.5,1\}$.
\end{exercise}

\begin{exercise}[Partial-rollout bias]
Using the simulator, record the mean length of samples consumed in each step under
\texttt{partial} and compare with the generated length distribution. How many steps does it
take for the consumed-length distribution to match the generated one? How does
\code{oversample} change this?
\end{exercise}

\begin{exercise}[Weight sync across sites]
A 32B policy is trained in one datacenter and rolled out in another over a 100\,Gb/s link.
Estimate the sync time in bf16, fp8 and with LoRA rank 64 on all linear layers. Which
schedule from \cref{sec:rlsys:async} would you use?
\end{exercise}

\begin{exercise}[Attacking the grader]
List five ways a policy could obtain reward from \cref{lst:rlsys:sandbox} without solving the
task, and give a mitigation for each. Implement two of them as tests for your own sandbox.
\end{exercise}

\begin{exercise}[Measurement]
Run \texttt{bench.py sim} and add a policy \texttt{async+partial} to the simulator (partial
rollouts in a generator pool that runs concurrently with training). Plot throughput and
mean staleness against \code{oversample}.
\end{exercise}

\begin{exercise}[Design]
You have 32 H100s and a 14B model. Propose a placement, schedule and parallel layout for an
RLVR run with $B=256$, $G=16$, mean response length 6{,}000 and a cap of 32{,}000 tokens.
Justify each choice with a napkin estimate.
\end{exercise}

% =====================================================================================
\section{Lab 29: Building blocks of an RL training system}\label{sec:rlsys:lab}

\begin{labbox}{29}{Building blocks of an RL training system}
\labdir{lab29_rlsys}
\labsection{Goals}
Implement and measure the pieces that surround the GRPO update in a production RL framework:
packed scoring passes, mismatch correction, schedule simulation and sandboxed rewards.

\labsection{Setup}
\begin{lstlisting}[style=shell]
cd code
pytest -q labs/lab29_rlsys                    # grades starter.py
LAB_IMPL=solution pytest -q labs/lab29_rlsys  # reference implementation
\end{lstlisting}
Everything runs on a CPU in under a minute with tiny models; \texttt{bench.py} accepts a real
Qwen3 checkpoint for the mismatch measurement.

\labsection{Tasks}
\begin{enumerate}
\item \textbf{Packing.} Implement \ident{ffd_pack} and \ident{efficiency_report} (padded,
length-sorted padded, packed real-token fractions).
\item \textbf{Varlen log-probs.} Implement \ident{block_causal_mask} and
\ident{packed_completion_logprobs} on top of the provided \ident{varlen_forward}. The tests
compare against unpacked \ident{sequence_logprob} for both Qwen3- and GPT-shaped models.
\item \textbf{Mismatch.} Implement \ident{is_weights} (truncate/mask/none),
\ident{mismatch_metrics} (including the $k_3$ KL estimate) and \ident{clipped_pg_loss} with an
optional TIS factor. Check that the gradient at $\pitheta=\piold$ is $-A/n_{\text{tokens}}$.
\item \textbf{Pipeline simulation.} Implement \ident{run_generation} (continuous batching with
an early stop) and \ident{weight_sync_seconds}; the provided \ident{simulate} drives your code
through the colocated, sync, async and partial policies.
\item \textbf{Sandboxed reward.} Implement \ident{run_python_reward} and \ident{score_batch}
with timeouts, resource limits and an empty environment.
\end{enumerate}

\labsection{Measurements}
\begin{itemize}
\item \texttt{bench.py sim}: throughput, utilisation, staleness and mixed-policy fraction for
each policy as $\sigma$ varies. Expect async to help most for light tails and partial rollouts
to dominate for heavy tails (\cref{fig:rlsys:simresults}).
\item \texttt{bench.py packing}: wall time of padded vs.\ packed log-prob passes and the
real-token fractions. Expect a large gap whenever lengths are heavy-tailed.
\item \texttt{bench.py mismatch}: bf16-vs-fp32 mismatch metrics, ideally on real Qwen3-0.6B
weights; look at the maximum ratio and the fraction outside $[0.5, 2]$.
\end{itemize}

\labsection{Deliverables}
Passing tests; a one-page note with the simulator sweep (a plot like
\cref{fig:rlsys:simresults}), the packing timings, and the mismatch metrics, each with a
sentence of interpretation.

\labsection{Stretch goals}
\begin{itemize}
\item Add an \texttt{async+partial} policy and a staleness-aware sampler that drops samples
older than $s$ versions; plot throughput vs.\ staleness.
\item Wire \ident{clipped_pg_loss} with TIS into your GRPO loop from \cref{ch:grpo2}, generate
rollouts with a bf16 copy of the model and train in fp32; compare reward curves with and without
TIS.
\item Replace the dense mask in \ident{varlen_forward} with a loop over segments that calls
\texttt{F.scaled\_dot\_product\_attention} per segment, and measure the speedup.
\end{itemize}
\end{labbox}
